\chapter{Run It and Verify the Result}

\begin{goalbox}
Start the bot, complete a real Telegram conversation, and run offline checks.
\end{goalbox}

Confirm that the environment is active and \texttt{BOT\_TOKEN} is set, then run:

\begin{lstlisting}[language=bash,numbers=none]
python bot.py
\end{lstlisting}

\begin{resultbox}
\texttt{Bot is running. Press Ctrl+C to stop it.}
\end{resultbox}

Leave that terminal open. In Telegram, open the link BotFather gave you or search
for the username you created. Press \textbf{Start}.

\section{A four-part manual check}

\begin{enumerate}
\item Send \texttt{/start}. The reply begins with your Telegram first name and
shows both commands.
\item Send \texttt{/help}. The reply gives three short usage steps.
\item Send \texttt{My first bot}. The reply must be
\texttt{You said: My first bot}.
\item Send a photo. The reply must be
\texttt{Please send a text message so I can repeat it.}
\end{enumerate}

If all four results match, the complete path---Telegram, your token, Python,
TeleBot, and the handlers---works.

\section{Run the offline tests}

Stop the bot with \texttt{Ctrl+C}, then run:

\begin{lstlisting}[language=bash,numbers=none]
python -m unittest -v
\end{lstlisting}

\begin{resultbox}
Four tests are listed, followed by \texttt{OK}. These checks use no token and do
not contact Telegram; they verify the reply-building functions.
\end{resultbox}

Whenever you change a message in \texttt{bot.py}, update its matching test and
run this command again.

\begin{checkpointbox}
Your bot answers all four manual checks, and the four offline tests pass.
\end{checkpointbox}
